\documentclass[compacto,10pt]{pubciencias-notas}

% -- Paquetes adicionales
\usepackage{comandosEPN}

% -- Datos del libro
\autor{Departamento de Formación Básica}
\materia{Álgebra Lineal}
\nota{Semana 05 Clase 02}
\fecha{Semestre 2025-A}

\logouno[1.1cm]{Logo_EPN.eps}
\logodos[1.7cm]{Logo_DFB.eps}

\definecolor{colordef}{cmyk}{0.81,0.62,0.00,0.22}

\begin{document}

\encabezado

\vspace*{-8mm}

\setcounter{section}{2}
\section{Determinantes}
\setcounter{subsection}{1}
\subsection{Propiedades fundamentales}

\begin{teo}[Determinante de una matriz triangular]
    Sea $A \in \Mat[\K]{n}{n}$ una matriz triangular superior o triangular inferior, entonces
    \[
        \det(A) = a_{11}a_{22} \cdots a_{nn},
    \]
    es decir, el determinante de una matriz triangular es el producto de los elementos de su diagonal principal.
\end{teo}

\begin{ejem}
    Consideremos la matriz
    $A = \begin{pmatrix}
    2 & 0 & 0 & 0 & 0 & 0 \\
    8 & -1 & 0 & 0 & 0 & 0 \\
    -7 & 6 & 2 & 0 & 0 & 0 \\
    4 & 0 & 3 & -2 & 0 & 0 \\
    -5 & 2 & 3 & 5 & 1 & 0 \\
    1 & -3 & 7 & -6 & 2 & 3
    \end{pmatrix}
    $.
    
    \noindent Se sigue que $\det(A) = (2)(-1)(2)(-2)(1)(3) = 24$
\end{ejem}

\begin{ejer}
    Calcular el determinante de la siguiente matriz
    \[
        B =
        \begin{pmatrix}
            -2 & 0 & 3 & 5 & -1 \\
            0 & 5 & -1 & 4 & 0 \\
            0 & 0 & -1 & 7 & -2 \\
            0 & 0 & 0 & 3 & -1 \\
            0 & 0 & 0 & 0 & 1
        \end{pmatrix}.
    \]
\end{ejer}

% \begin{prop}
%     Sea $E\in\Mat[\K]{n}{n}$ una matriz elemental, $i,j\in I$ y $\alpha\in\K$, con $\alpha\neq 0$. Se tiene que
%     \begin{itemize}
%     \item   
%         si la operación es $F_i \leftrightarrow F_j$, entonces $\det(E) = - 1$;
%     \item 
%         si la operación es $\alpha F_i \rightarrow F_i$, entonces $\det(E) = \alpha$; y
%     \item 
%         si la operación es $\alpha F_i + F_j\rightarrow F_j$, entonces $\det(E) = 1$.
%     \end{itemize}
% \end{prop}

% \begin{ejem}
%     Consideremos la matriz $
%     \begin{pmatrix}
%         1 & 0 \\
%         0 & 1
%     \end{pmatrix}$, cuyo determinante es $1$.
    
%     \begin{itemize}
%     \item
%         Si hacemos $F_{1} \leftrightarrow F_{2}$,  obtenemos $ E =
%         \begin{pmatrix}
%             0 & 1 \\
%             1 & 0
%         \end{pmatrix}$. Se sigue que $\det(E) = -1$.
        
%     \item
%         Si hacemos $3F_{1} \rightarrow F_{1}$, obtenemos $ E =
%         \begin{pmatrix}
%             3 & 0 \\
%             0 & 1
%         \end{pmatrix}$. Se sigue que $\det(E) = 3$.
        
%     \item
%         Si hacemos $-2F_{2} + F_{1} \rightarrow F_{1} $, obtenemos $ E =
%         \begin{pmatrix}
%             1 & -2 \\
%             0 & 1
%         \end{pmatrix}$.
        
%         \noindent Se sigue que $\det(E) = 1$.
%     \end{itemize}
% \end{ejem}

% \begin{ejer}
%     Considere la matriz $ 
%     \begin{pmatrix}
%         1 & 0 & 0 \\
%         0 & 1 & 0 \\
%         0 & 0 & 1
%     \end{pmatrix}$, cuyo determinante es $1$. Para cada operación por fila determine a la correspondiente matriz elemental y a su determinante.
%     \begin{enumerate}
%         \item[$a)$] $F_{3} \leftrightarrow F_{1}$
%         \item[$b)$] $-2F_{2} \rightarrow F_{2}$
%         \item[$c)$] $3F_{3} + F2 \rightarrow F_{2}$
%     \end{enumerate}
% \end{ejer}

\begin{prop}
    Sea $A \in \Mat[\K]{n}{n}$. El determinante de la matriz $A$ y de su transpuesta son iguales, es decir, 
    \[
        \det(A^\intercal) = \det(A).
    \]
\end{prop}

\begin{obs}
    Por esta proposición, cuando se \textbf{calcula el valor de un determinante}, da lo mismo trabajar por filas que trabajar por columnas.
\end{obs}

\begin{teo}[Determinante del producto de matrices]\label{DetProd}
    Sean $A, B\in \Mat[\K]{n}{n}$, se tiene que:
    \[
        \det(AB) = \det(A)\det(B).
    \]
    es decir, el determinante de un producto de matrices es igual al producto de los determinantes de las matrices.
\end{teo}

\begin{ejem} Dadas las matrices 
\[
A = 
    \begin{pmatrix}
        -2 & 3 \\
        1 & 1
    \end{pmatrix}
\texty
B=
    \begin{pmatrix}
        1 & -1 \\
        2 & 4
    \end{pmatrix}.
\]
Se sigue que $AB = 
    \begin{pmatrix}
        4 & 14 \\
        3 & 3
    \end{pmatrix}$ y ya que $\det(A) = -5$, $\det(B) = 6$ y $\det(AB) = -30$, la igualdad del teorema \ref{DetProd} se verifica.
\end{ejem}

\begin{ejer}\hspace{0pt} Dadas las matrices
\[
A = 
    \begin{pmatrix}
        3 & 0 & 3 \\
        0 & 3 & 0 \\
        3 & 0 & 3
    \end{pmatrix}
\texty    
B = 
    \begin{pmatrix}
        -2 & 0 & -2 \\
        0 & -2 & 0 \\
        2 & 0 & 2
    \end{pmatrix}.
\]
Calcular $AB$, $\det(AB)$, $\det(A)$, $\det(B)$ y compruebe la igualdad del teorema \ref{DetProd}.
\end{ejer}

\begin{prop}
    Sea $A \in \Mat[\K]{n}{n}$. Si una fila o columna de $A$ contiene solo ceros, entonces $\det (A) =0$.
\end{prop}

\begin{ejem}
    Sea $A =
    \begin{pmatrix}
        3 & -1 & 11 & 0 \\
        9 & 7 & -5 & 0 \\
        -2 & -6 & -8 & 0 \\
        4 & 1 & -5 & 0
    \end{pmatrix}$. Realizamos el cálculo del determinante de $A$ mediante la expansión por cofactores de $A$, con respecto a la columna $4$. Obtenemos,
    \begin{align*}
        \det(A) &= \sum_{k=1}^{4} a_{k4} C_{k4}\\
                &= \sum_{k=1}^{4} 0\cdot C_{k4}\\
                &= 0
    \end{align*}
    
\end{ejem}

\begin{prop}
    Sea $A \in \Mat[\K]{n}{n}$. Si dos filas o columnas de $A$ son iguales, entonces 
    \[
        \det(A) = 0
    \]
\end{prop}

\begin{ejem}
    Sea $A = 
    \begin{pmatrix}
        1 & 3 & 1 & 2 \\
        -1 & 1 & -1 & 3 \\
        2 & 3 & 2 & 1 \\
        -2 & 1 & -2 & 3
    \end{pmatrix}$, calculemos su determinante mediante la expansión por cofactores con respecto a la primera columna.
    \begin{align*}
        \det(A) = & 1
        \begin{vmatrix}
            1 & -1 & 3 \\
            3 & 2 & 1 \\
            1 & -2 & 3
        \end{vmatrix} + 1
        \begin{vmatrix}
            3 & 1 & 2 \\
            3 & 2 & 1 \\
            1 & -2 & 3
        \end{vmatrix} +2
        \begin{vmatrix}
            3 & 1 & 2 \\
            1 & -1 & 3 \\
            1 & -2 & 3
        \end{vmatrix} +2
        \begin{vmatrix}
            3 & 1 & 2 \\
            1 & -1 & 3 \\
            3 & 2 & 1
        \end{vmatrix} \\
        = & -8 + 0 + 2(7) + 2(-3) \\
        = & \, 0.
    \end{align*}
\end{ejem}

\begin{ejer} Calcular el determinante de la matriz $B = 
    \begin{pmatrix}
        1 & -1 & 1 \\
        0 & 5 & 0 \\
        2 & 3 & 2
    \end{pmatrix}$.
\end{ejer}

\begin{prop}
    Sea $A \in \Mat[\K]{n}{n}$. Si una  fila o columna de $A$ es múltiplo de otra, entonces 
    \[
        \det(A) = 0
    \]
\end{prop}

\begin{ejem}
    El determinante de la matriz $A =
\begin{pmatrix}1 & 2 & -1 \\-2 & 5 & -2 \\3 & 6 & -3\end{pmatrix}$ es nulo. Pues la fila $3$ es múltiplo de la fila $1$. Por tanto, $\det(A)=0$
\end{ejem}

\begin{ejer} Calcular el determinante de la matriz $B = 
    \begin{pmatrix}
        2 & -6 & 4 \\
        -1 & 3 & 1 \\
        3 & -9 & 3
    \end{pmatrix}$.
\end{ejer}

% \begin{ejer}
%     Verifique que el determinante de la matriz $B$ del ejemplo $12$ es cero, realizando el cálculo de otra manera.
% \end{ejer}

% \begin{ejer} Demuestre a la proposición $3$.
% \end{ejer}

\begin{prop}\label{DetSign}
    Sean $A, B \in \Mat[\K]{n}{n}$. Si la matriz $B$ se obtiene intercambiando dos filas o columnas de $A$ entonces
    \[  
        \det(B) = - \det(A).
    \]
\end{prop}

\begin{ejem} Dadas las matrices 
\[
A = 
    \begin{pmatrix}
        1 & 0 & 0 \\
        3 & -1 & 0 \\
        3 & 3 & -2 
    \end{pmatrix}
\texty
B =
    \begin{pmatrix}
        1 & 0 & 0 \\
        3 & 3 & -2 \\
        3 & -1 & 0
    \end{pmatrix}.
\]
Notemos que la matriz $B$ se obtiene al intercambiar las filas $2$ y $3$ de la matriz $A$. Además, $\det(A) = 2$ y $\det(B) = -2$, es decir, la igualdad de la proposición \ref{DetSign} se verifica.
\end{ejem}

\begin{ejer}\hspace{0pt} Dadas las matrices
\[
A = 
    \begin{pmatrix}
        1 & 0 & 0 \\
        3 & -1 & 0 \\
        3 & 3 & -2 
    \end{pmatrix}
\texty
B =
    \begin{pmatrix}
        -1 & 3 & 0 \\
        0 & 1 & 0 \\
        3 & 3 & -2
    \end{pmatrix}.
\]
Comparar a sus determinantes.
\end{ejer}

\begin{prop}
    Sean $A, B \in \Mat[\K]{n}{n}$. Si $B$ se obtiene al multiplicar una fila o columna de $A$ por un escalar $\alpha \in \K$, entonces
    \[
        \det(B) = \alpha \det(A).
    \]
\end{prop}

\begin{ejem}\hspace{0pt} Dada la matriz $A =
\begin{pmatrix}1 & -1 & -1 \\-2 & 1 & -2 \\0 & 2 & 0\end{pmatrix}$ cuyo determinante es 8. Considere a la matriz $B =
\begin{pmatrix}1 & 2 & -1 \\-2 & -2 & -2 \\0 & -4 & 0\end{pmatrix}$. Notemos que $B$ se obtiene al multiplicar la segunda columna de la matriz $A$ por $-2$.
    
    \noindent Luego, $\det(B) = -2(8) = -16$.
\end{ejem}

% \begin{ejer} Demuestre a la proposición $9$.
% \end{ejer}

\begin{prop}
    Sea $A\in \Mat[\K]{n}{n}$ y $\alpha \in\K$. Se tiene que
    \[
        \det(\alpha A) = \alpha^n \det(A).
    \]
\end{prop}

\begin{ejem}
    Sea $A =
    \begin{pmatrix}
        -2 & -3 & 6 & 4 \\
        0 & 2 & 6 & -9 \\
        0 & 0 & -1 & 8 \\
        0 & 0 & 0 & 2
    \end{pmatrix}$. Notemos que al ser $A$ una matriz triangular superior, $\det(A) = 8$.
    
    \noindent Consideremos la matriz $-2A = 
    \begin{pmatrix}
        4 & 6 & -12 & -8 \\
        0 & -4 & -12 & 18 \\
        0 & 0 & 2 & -16 \\
        0 & 0 & 0 & -4
    \end{pmatrix}$ y calculemos su determinante. Puesto que está también es una matriz triangular superior, se sigue que $\det(-2A) = 128 = (-2)^{4}8$.
\end{ejem}

\begin{ejer}
    Sea $A = 
    \begin{pmatrix}
        1 & -2 \\
        -1 & 3
    \end{pmatrix}$. ¿Cuál es el valor de $\det(-4A)$?
\end{ejer}

\begin{prop}
    Sean $A, B \in \Mat[\K]{n}{n}$, $\alpha\in\K$ e $i,j\in I$, con $i\neq j$. Si $B$ se obtiene al aplicar una operación de fila $\alpha F_i + F_j \to F_j$, entonces 
    \[
        \det(B) = \det(A).
    \]
\end{prop}

% \begin{prop} Sean $A, B, C \in \Mat[\K]{n}{n}$ y $j\in I$ tales que $A$, $B$ y $C$ difieran únicamente en la columna $j$. Si la columna $j$ de $C$ es el resultado de sumar las columnas $j$ de $A$ y $B$, entonces 
%     \[
%         \det(C) = \det(A) + \det(B).
%     \]
% \end{prop}

% \begin{ejem}
%     Si $\det
%         \begin{pmatrix}
%             x&y&z\\
%             3&0&2\\
%             1&1&1
%         \end{pmatrix}
%     =1$, calcular el determinante de la matriz:
%     \[
%         \begin{pmatrix}
%             x&y&z\\
%             3x+3&3y&3z+2\\
%             x+1&y+1&z+1
%         \end{pmatrix}
%     \]
    % Notemos que
    % \begin{align*}
    %     \det\begin{pmatrix}
    %                 x & y & z\\
    %                 3x+3 & 3y & 3z+2\\
    %                 x+1 & y+1 & z+1
    %             \end{pmatrix}
    %         &= 
    %         \det\begin{pmatrix}
    %                 x&y&z\\
    %                 3x & 3y & 3z\\
    %                 x+1 & y+1 & z+1
    %             \end{pmatrix}\\
    %         &\phantom{=} 
    %         + \det
    %             \begin{pmatrix}
    %                 x&y&z\\
    %                 3 & 0 & 2\\
    %                 x+1 & y+1 & z+1
    %             \end{pmatrix} \\
    %         &= 
    %         3\det\begin{pmatrix}
    %                 x&y&z\\
    %                 x & y & z\\
    %                 x+1 & y+1 & z+1
    %             \end{pmatrix}\\
    %         &\phantom{=} + \det
    %             \begin{pmatrix}
    %                 x&y&z\\
    %                 3 & 0 & 2\\
    %                 x & y & z
    %             \end{pmatrix}\\
    %         &\phantom{=} + \det
    %             \begin{pmatrix}
    %                 x&y&z\\
    %                 3 & 0 & 2\\
    %                 1 & 1 & 1
    %             \end{pmatrix}\\
    %         &= 3(0) + 0 + 1 \\
    %         &= 1. 
    % \end{align*}
    % Por lo tanto,
    % \[
    %     \det\begin{pmatrix}
    %         x & y & z\\
    %         3x+3 & 3y & 3z+2\\
    %         x+1 & y+1 & z+1
    %     \end{pmatrix} = 1.
    % \]
% \end{ejem}

\begin{ejer}
    Si $\det
        \begin{pmatrix}
            x&y&z\\
            3&0&2\\
            1&1&1
        \end{pmatrix}
    =1$, calcular el determinante de la matriz:
    
    \[
        \begin{pmatrix}
            x & y & z \\
            3x+3 & 3y & 3z+2 \\
            2 & 2 & 2
        \end{pmatrix}.
    \]
\end{ejer}

\begin{ejer}
  Utilizando propiedades de determinantes, calcular:
   \[
    \begin{vmatrix}
            1 & a & a^2 \\
            1 & b & b^2 \\
            1 & c & c^2
        \end{vmatrix}
    \]
\end{ejer}
%\begin{prop} Sean $A, B, C \in \Mat[\K]{n}{n}$ y $j\in I$ tales que $A$, $B$ y $C$ difieran únicamente en la columna $j$. Si la columna $j$ de $C$ es el resultado de sumar las columnas $j$ de $A$ y $B$, entonces 
 %   \[
  %      \det(C) = \det(A) + \det(B).
   % \]
%\end{prop}

%\begin{ejem}
 %   Si $\det\begin{pmatrix}
  %      -2 & -3 & 1 \\
   %     0 & 2 & x\\
    %    1 & 1 & y
    %\end{pmatrix} = 8$ y %$\det\begin{pmatrix}
     %   -2 & 6 & 1 \\
      %  0 & 4 & x\\
       % 1 & -3 & y
    %\end{pmatrix} = 5$.
    
    %Entonces, 
    %\begin{align*}
     %   \det\begin{pmatrix}
      %  -2 & 3 & 1 \\
       % 0 & 6 & x\\
       % 1 & -2 & y
    %\end{pmatrix} &= \det\begin{pmatrix}
    %    -2 & -3 & 1 \\
     %   0 & 2 & x\\
      %  1 & 1 & y
  %  \end{pmatrix} + \det\begin{pmatrix}
   %     -2 & -3 & 1 \\
    %    0 & 2 & x\\
     %   1 & 1 & y
    %\end{pmatrix}\\
    %&= 8 + 5 \\
    %&= 13
    %\end{align*}

%\end{ejem}

\end{document}